\begin{textsample}{POS Dim 1 – vem\_ed\_23 – Score 5.00 – vem\_ed\_23/t119.txt}  \label{ex:f1_pos_145}
Conferência Faubai 2024: \textbf{Ações} para a internacionalização Como moderador, Succi Junior atuou na sessão sobre Intercâmbios Virtuais agendada para 23 de abril, das 9h às 11h.
E participou também da mesa-redonda “Intercâmbio Virtual: \textbf{ações} realizadas na UEM, Unesp e Centro Paula Souza”, das 17h às 18h30 no dia 23 de abril.
Ana Carolina Freschi, responsável pelos PCIs em inglês, discorreu sobre o \textbf{papel} dos professores de línguas estrangeiras como promotores da internacionalização. “Virtual Exchange in Technological Higher Education: The Role of Foreign Language Teachers as Promoters of Internationalization” \textbf{foi} um trabalho elaborado em coautoria com Neusa Haruka Sezaki Gritti, responsável pelo \textbf{apoio} pedagógico e administrativo aos PCIs / Cesu.
História da Conferência Faubai Desde 1988, a Associação Brasileira de Educação Superior (Faubai) organiza eventos anuais sobre internacionalização: a Conferência Faubai \textbf{é} o maior e mais antigo evento sobre o tema na América Latina.
Em 2024, mais de 700 participantes de 25 países se reuniram no Centro de Convenções Rebouças, na capital paulista, para discutir a diversidade de \textbf{contextos} da internacionalização plural, por meio de abordagens críticas, antirracistas e sintonizadas com as novas tecnologias da informação e da comunicação."Temos o importante desafio de buscar uma internacionalização \textbf{que} entenda a cooperação internacional a \textbf{partir} de novas óticas, \textbf{cada} vez mais inclusivas e diversas, e incorpore novas visões e vozes de \textbf{todas} as regiões do mundo, novos significados, novas epistemologias", ressalta José Celso Freire Junior, Presidente da FAUBAI e Assessor de Relações Externas da Unesp.
Afinal, o tema de 2024 \textbf{é} “Inviting for a new Journey” (em tradução livre, “Convite para uma nova jornada”).
O público-alvo do evento compreende gestores de relações internacionais de instituições de educação \textbf{superior} (IES); reitores, vice-reitores e pró-reitores de IES; especialistas em educação internacional; representantes de organizações governamentais e não-governamentais; representantes de organismos diplomáticos; membros de organizações empresariais; docentes e pesquisadores.

% matched lemmas: apoio, ação, cada, contexto, papel, partir, que, ser, superior, todo
\end{textsample}
